\section{Introduction}

Standard residual connections~\cite{he2015resnet} are the \emph{de facto} building block of modern LLMs~\cite{openai-2024-gpt4,touvron-2023-llama,deepseekaiv3}. The update $\bm{h}_l = \bm{h}_{l-1} + f_{l-1}(\bm{h}_{l-1})$ is widely understood as a \textit{gradient highway} that lets gradients bypass transformations via identity mappings, enabling stable training at depth. Yet residuals also play a second role that has received less attention. Unrolling the recurrence shows that every layer receives the same uniformly-weighted sum of all prior layer outputs; residuals define how information aggregates across depth. Unlike sequence mixing and expert routing, which now employ learnable input-dependent weighting~\cite{transformer,adaptive,deepseekaiv3}, this depth-wise aggregation remains governed by fixed unit weights, with no mechanism to selectively emphasize or suppress individual layer contributions.


In practice, PreNorm~\cite{xiong2020layer} has become the dominant paradigm, yet its unweighted accumulation causes hidden-state magnitudes to grow as $O(L)$ with depth, progressively diluting each layer's relative contribution~\cite{li2026siamesenorm}. Early-layer information is buried and cannot be selectively retrieved; empirically, a significant fraction of layers can be pruned with minimal loss~\cite{gromov2025unreasonable}. Recent efforts such as scaled residual paths~\cite{wang2022deepnorm} and multi-stream recurrences~\cite{zhu2025hyperconnections} remain bound to the additive recurrence, while methods that do introduce cross-layer access~\cite{pagliardini2024denseformer,xiao2025muddformer} are difficult to scale. The situation parallels the challenges that recurrent neural networks (RNNs) faced over the sequence dimension before attention mechanism provided an alternative.

We observe a formal duality between depth-wise accumulation and the sequential recurrence in RNNs. Building on this duality, we propose \textbf{Attention Residuals (AttnRes)}, which replaces the fixed accumulation $\bm{h}_l = \sum_{i} \bm{v}_i$ with $\bm{h}_l = \sum_{i} \brickred{\alpha_{i \to l}} \cdot \bm{v}_i$, where $\brickred{\alpha_{i \to l}}$ are $\operatorname{softmax}$ attention weights computed from a single learned pseudo-query $\bm{w}_l \in \mathbb{R}^d$ per layer. This lightweight mechanism enables selective, content-aware retrieval across depth with only one $d$-dimensional vector per layer. Indeed, standard residual connections and prior recurrence-based variants can all be shown to perform depth-wise \emph{linear} attention; AttnRes generalizes them to depth-wise $\operatorname{softmax}$ attention, completing for depth the same linear-to-$\operatorname{softmax}$ transition that proved transformative over sequences (\S\ref{sec:structured-depth}, \S\ref{sec:depth-ttt}).

In standard training, Full AttnRes adds negligible overhead, since the layer outputs it requires are already retained for backpropagation. At scale, however, activation recomputation and pipeline parallelism are routinely employed, and these activations must now be explicitly preserved and communicated across pipeline stages. We introduce \textit{Block AttnRes} to maintain efficiency in this regime: layers are partitioned into $N$ blocks, each reduced to a single representation via standard residuals, with cross-block attention applied only over the $N$ block-level summaries. This brings both memory and communication down to $O(Nd)$, and together with infrastructure optimizations (\S\ref{sec:infra}), Block AttnRes serves as a drop-in replacement for standard residual connections with marginal training cost and negligible inference latency overhead.

Scaling law experiments confirm that AttnRes consistently outperforms the baseline across compute budgets, with Block AttnRes matching the loss of a baseline trained with $1.25\times$ more compute. We further integrate AttnRes into the Kimi Linear architecture~\cite{zhang2025kimi} (48B total / 3B activated parameters) and pre-train on 1.4T tokens. Analysis of the resulting training dynamics reveals that AttnRes mitigates PreNorm dilution, with output magnitudes remaining bounded across depth and gradient norms distributing more uniformly across layers. On downstream benchmarks, our final model improves over the baseline across all evaluated tasks.

\textbf{Contributions}

\begin{itemize}[leftmargin=*]
    \item \textbf{Attention Residuals.} We propose AttnRes, which replaces fixed residual accumulation with learned $\operatorname{softmax}$ attention over depth, and its scalable variant Block AttnRes that reduces memory and communication from $O(Ld)$ to $O(Nd)$. Through a unified structured-matrix analysis, we show that standard residuals and prior recurrence-based variants correspond to depth-wise \emph{linear} attention, while AttnRes performs depth-wise $\operatorname{softmax}$ attention.
    \item \textbf{Infrastructure for scale.} We develop system optimizations that make Block AttnRes practical and efficient at scale, including cross-stage caching that eliminates redundant transfers under pipeline parallelism and a two-phase inference strategy that amortizes cross-block attention via online $\operatorname{softmax}$~\citep{milakov2018online}. The resulting training overhead is marginal, and the inference latency overhead is less than 2\% on typical inference workloads.
    \item \textbf{Comprehensive evaluation and analysis.} We validate AttnRes through scaling law experiments, component ablations, and downstream benchmarks on a 48B-parameter model pre-trained on 1.4T tokens, demonstrating consistent improvements over standard residual connections. Training dynamics analysis further reveals that AttnRes mitigates PreNorm dilution, yielding bounded hidden-state magnitudes and more uniform gradient distribution across depth.
\end{itemize}

\newpage